% P18-A1: the exact Dung Section 4.3 constructions and corrected bridges.
\chapter{Logic Programs as Argumentation Frameworks}
\label{sec:p18-logic-program-argumentation}

Dung gives two different argument constructions for normal logic
programs. One uses finite derivations supported by negative assumptions;
the other uses individual ground clauses. Their grounded semantics
give respectively the well-founded and Fitting models. Their stable
extensions give respectively stable models and models of completion
\cite[sec.~4.3 and appendices~A--C]{Dung1995}.

\begin{hbdefinition}{def:p18-program-syntax}{Programs, ground literals and interpretations}
A normal clause has the form
\[
 p\leftarrow q_1,\ldots,q_m,\neg r_1,\ldots,\neg r_n
                 \qquad(m,n\geq0),
\]
with an atomic head and a finite body. The symbol \(\neg\) in
a body denotes default negation; no negative heads are admitted.
A source program is a finite set of such first-order clauses.
Use its free term algebra, adding a fresh constant if necessary
to make the Herbrand universe nonempty. Let \(H\) be its
Herbrand base and \(Q=G_P\) its possibly infinite set of ground
instances. The semantic constructions below also work for any
set \(Q\) of ground clauses over a specified set \(H\), provided
every body is finite. Claims about syntactic decidability and
the finite predicate dependency graph retain the finite-source
program hypothesis.

Let \(\operatorname{Lit}(H)=H\cup\neg H\), where
\(\neg X=\{\neg p:p\in X\}\), and let \(l^*\) be the
complement of a literal. A partial interpretation is a
consistent subset \(I\subseteq\operatorname{Lit}(H)\).
Write \(I^+=I\cap H\) and
\(I^-=\{p:\neg p\in I\}\); these superscripts in this
chapter are truth and falsity components, not attack targets.
A finite body \(K\) is true in \(I\) when \(K\subseteq I\),
false when \(l^*\in I\) for some \(l\in K\), and otherwise
undefined. The empty body is true and is never false.
A total interpretation is
\(\widehat M=M\cup\neg(H\setminus M)\), for \(M\subseteq H\).
\end{hbdefinition}

\begin{hbdefinition}{def:p18-program-reduct}{Positive consequence, finite support and the reduct}
For a positive ground program \(D\), put
\[
 T_D(X)=\{p:\exists(p\leftarrow K)\in D,\ K\subseteq X\},
 \qquad \operatorname{LM}(D)=\bigcup_{n<\omega}T_D^n(\varnothing).
\]
For finite \(K\subseteq\neg H\), \(K\) supports \(p\)
when a finite derivation of \(p\) uses clauses of \(Q\),
derives every positive body atom earlier, and uses only
negative body literals in \(K\). Equivalently a finite
proof tree has root \(p\), expands a positive node by
a ground clause, and has only negative assumption leaves
and fact leaves. Every body occurrence is expanded.
The support need not be minimal and may contain unused
negative assumptions.

The reduct \(Q^X\) deletes clauses having a negative
body literal \(\neg p\) with \(p\in X\), then deletes
the remaining negative literals. Set
\(\Gamma(X)=\operatorname{LM}(Q^X)\).
A stable model is \(M=\Gamma(M)\).
\end{hbdefinition}

\begin{hblemma}{thm:p18-support-reduct}{proved=annex:p18.support.reduct}{Finite proof witnesses and the exact reduct interface}
\(\operatorname{LM}(D)\) is the least model and least
fixed point of \(T_D\). An atom is in it exactly when
it has a finite proof tree. For every \(X\subseteq H\),
\[
 p\in\Gamma(X)\quad\Longleftrightarrow\quad
 \text{some finite support }K\subseteq\neg(H\setminus X)
                                      \text{ supports }p.
\]
The map \(\Gamma\) reverses inclusion. All claims allow
infinitely many ground clauses and infinitely many
clauses with the same head.
\end{hblemma}

\begin{hbdefinition}{def:p18-source-program-graphs}{The two literal source argument graphs}
The support graph \(\mathcal A_\infty(Q)\), Dung's
\(AF_{\rm napif}(P)\), has arguments
\[
 \{(K,p):K\subseteq\neg H\text{ finite and supports }p\}
       \ \cup\ \{\nu_p:p\in H\},\qquad
                        \nu_p=(\{\neg p\},\neg p).
\]
The clause graph \(\mathcal A_{\rm c}(Q)\), Dung's
\(AF_{\rm naff}(P)\), instead has
\[
 \{(K,p):(p\leftarrow K)\in Q\}\ \cup\ \{\nu_p:p\in H\}.
\]
Here \(K\) can contain both positive and negative literals.
Membership in either graph below means membership
in its argument set.
Repeated body occurrences are identified when taking its
set; finite conjunction is unchanged.
In both graphs the exact attack clause is
\[
                  (K,h)\to(K',h')\quad\Longleftrightarrow\quad h^*\in K'.
\]
Write \(\operatorname{con}(E)=\{h:\exists K\ (K,h)\in E\}\).
Use \(F_\infty,F_{\rm c}\) for the two characteristic
operators, and \(N_\infty(E)\) for the arguments of
\(\mathcal A_\infty\) not attacked by \(E\).
For \(B\subseteq H\), let
\[
 {\cal E}_\infty(B)=
   \{(K,p):K\subseteq\neg B\text{ and }(K,p)\in\mathcal A_\infty\}
                    \cup\{\nu_p:p\in B\}.
\]
\end{hbdefinition}

\begin{hblemma}{thm:p18-support-operator}{proved=annex:p18.support.operator}{The full support-graph operator identities}
For every argument set \(E\), put
\(T=\operatorname{con}(E)\cap H\). Then
\[
\begin{aligned}
 N_\infty(E)&={\cal E}_\infty(H\setminus T),\\
 F_\infty(E)&={\cal E}_\infty(H\setminus\Gamma(T)),\\
 \operatorname{con}(F_\infty(E))
             &=\Gamma^2(T)\cup\neg(H\setminus\Gamma(T)).
\end{aligned}
\]
If \(E\) is conflict-free its conclusion set is consistent.
If it is admissible, additionally
\[
 T\subseteq\Gamma^2(T),\qquad
       \operatorname{con}(E)^-\subseteq H\setminus\Gamma(T).
\]
\end{hblemma}

\begin{hbtheorem}{thm:p18-program-stable-correspondence}{proved=annex:p18.program.stable.correspondence}{The exact stable-model bijection}
Stable models of \(Q\) correspond bijectively to stable
extensions of the literal support graph:
\[
          M\longmapsto{\cal E}_\infty(H\setminus M).
\]
The inverse takes the positive conclusions; the full
conclusion set of the associated extension is \(\widehat M\).
This proves both directions of Dung's Theorem~49.
\end{hbtheorem}

\section{Unfounded sets and the well-founded model}

\begin{hbdefinition}{def:p18-unfounded-operators}{The exact well-founded operators}
For \(I\subseteq\operatorname{Lit}(H)\), let
\[
 T_Q(I)=\{p:\exists(p\leftarrow K)\in Q,\ K\subseteq I\}.
\]
A set \(U\subseteq H\) is unfounded with respect to \(I\)
when every clause with head in \(U\) has a false body in
\(I\), or has a positive body atom in \(U\).
Let \(U_Q(I)\) be the union of all such sets.
Define
\[
 V_Q(I)=T_Q(I)\cup\neg U_Q(I).
\]
For positive-consequence saturation, put
\[
 J_0=I,\quad J_{n+1}=J_n\cup T_Q(J_n),\quad
 T_Q^*(I)=\bigcup_{n<\omega}J_n,
 \qquad W_Q(I)=T_Q^*(I\cup\neg U_Q(I)).
\]
These formulas are also defined on possibly inconsistent
literal sets as an auxiliary powerset domain. The
least fixed points constructed below are consistent.
The well-founded model is \(\operatorname{lfp}(V_Q)\).
\end{hbdefinition}

\begin{hblemma}{thm:p18-unfounded-support}{proved=annex:p18.unfounded.support}{The greatest unfounded set and the full Appendix B interface}
The union \(U_Q(I)\) is itself unfounded and is the
greatest unfounded set. The operators \(U_Q,V_Q,W_Q\)
are monotone in information inclusion.
If \(E\) is admissible in \(\mathcal A_\infty\),
\(I=\operatorname{con}(E)\), and \(T=I^+\), then
\[
 U_Q(I)=H\setminus\Gamma(T)
       =\{p:\nu_p\in F_\infty(E)\},
 \qquad
        W_Q(I)=\operatorname{con}(F_\infty(E)).
\]
The proof gives both unfounded-set directions with
finite proof trees and the complete positive-saturation
step, including facts and empty supports.
\end{hblemma}

\begin{hbtheorem}{thm:p18-program-well-founded}{proved=annex:p18.program.well.founded}{Well-founded semantics and the literal support graph}
The least fixed points of \(V_Q\) and \(W_Q\) exist
and agree. Their common value is consistent and is
\[
            \operatorname{WFM}(Q)
                 =\operatorname{con}(G_{\mathcal A_\infty(Q)}).
\]
At every ordinal stage, iteration of \(W_Q\) from
the empty set equals the conclusions of the
corresponding grounded argument stage; both take
unions at nonzero limits. No finitary-graph
hypothesis is imposed. This supplies Theorem~50
and its Appendix~B construction.
\end{hbtheorem}

\section{Completion and the clause graph}

\begin{hbdefinition}{def:p18-completion-semantics}{Clark completion and the precise three-valued convention}
For each source clause \(p(\mathbf t)\leftarrow K\),
choose fresh \(\mathbf x\) and replace its body by
\[
        B(\mathbf x)=\exists\mathbf y\,
                 (\mathbf x=\mathbf t\land\bigwedge K),
\]
where \(\mathbf y\) lists all original variables.
The completed definition of \(p\) is
\[
         \forall\mathbf x\,
            \left(p(\mathbf x)\leftrightarrow\bigvee_{C\text{ with head }p}
                                                  B_C(\mathbf x)\right).
\]
An empty disjunction is false. Clark equality uses
identity of ground terms, distinct outer constructors,
injectivity of constructors, and exclusion of a
finite term equal to a proper term containing it,
together with equality congruence. These hold in
the free Herbrand term algebra used here.
For arbitrary ground \(Q\), the corresponding
supported-model condition is the pointwise equation
\[
             p\in M\quad\Longleftrightarrow\quad
                    \exists(p\leftarrow K)\in Q\ (K\subseteq\widehat M).
\]
An arbitrary infinite disjunction is not claimed
to be one finitary first-order formula.

For partial models, conjunction is true when all
conjuncts are true and false when one is false;
disjunction is true when one disjunct is true
and false when all are false; negation exchanges
true and false and preserves undefined.
Existential quantification has the analogous
some-true/all-false clauses. The completed head
and completed body must have the \emph{same}
value, including the case where both are undefined.
This is Fitting's truth-value equality connective,
not an ordinary strong-Kleene biconditional
\cite[sec.~5, pp.~299--301]{Fitting1985KripkeKleene}.
The exact ground operator is
\[
 Z_Q(I)=\{p:\text{every }p\text{-clause has a false body in }I\},
             \qquad\Phi_Q(I)=T_Q(I)\cup\neg Z_Q(I).
\]
A partial completion model is a consistent fixed
point of \(\Phi_Q\). The Fitting model is its
least fixed point.
\end{hbdefinition}

\begin{hbtheorem}{thm:p18-clause-completion}{proved=annex:p18.clause.completion}{Complete, preferred and stable clause-graph correspondences}
Define
\[
 {\cal E}_{\rm c}(I)=
 \{(K,p)\in\mathcal A_{\rm c}:K\subseteq I,\ p\in H\}
                          \cup\{\nu_p:p\in I^-\}.
\]
The map \(I\mapsto{\cal E}_{\rm c}(I)\) is an
order isomorphism from partial completion models
to complete extensions of \(\mathcal A_{\rm c}\),
with inverse \(\operatorname{con}\).
It sends maximal partial completion models to
preferred extensions and total completion models
to stable extensions, in both directions.
In particular the latter is the full bijection
of Theorem~55 and Appendix~C.
\end{hbtheorem}

\begin{hbtheorem}{thm:p18-clause-fitting}{proved=annex:p18.clause.fitting}{The correct accelerated operator and Fitting correspondence}
Put
\[
 \Psi_Q(I)=T_Q(I^+\cup\neg Z_Q(I))\cup\neg Z_Q(I).
\]
For every set \(E\) of clause-graph arguments,
\[
          \operatorname{con}(F_{\rm c}(E))
                         =\Psi_Q(\operatorname{con}(E)).
\]
The monotone operators \(\Psi_Q,\Phi_Q\) have
exactly the same fixed points on the auxiliary
literal powerset, and hence the same least
fixed point. It is consistent and equals
\[
       \operatorname{FM}(Q)
                         =\operatorname{con}(G_{\mathcal A_{\rm c}(Q)}).
\]
The equality is between the least fixed points.
The separate \(\Phi_Q\) and argument iterations
can take different successor steps.
\end{hbtheorem}

\section{Used supports and the corrected graph transfer}

\begin{hbdefinition}{def:p18-used-support-graph}{The explicitly distinct used-support reduction}
A used support of \(p\) is exactly the set
of negative leaves of some finite proof tree
of \(p\). Let \(\mathcal A_{\rm u}(Q)\)
retain only these positive arguments from
\(\mathcal A_\infty(Q)\), together with every
\(\nu_p\), and retain the same attack clause.
The literal source graph \(\mathcal A_\infty\)
is unchanged. Every support argument \((K,p)\)
has a used-support witness \((J,p)\) with \(J\subseteq K\).
\end{hbdefinition}

\begin{hbtheorem}{thm:p18-used-support-reduction}{proved=annex:p18.used.support.reduction}{The full semantic bridge for support reduction}
The complete-extension posets of
\(\mathcal A_\infty\) and \(\mathcal A_{\rm u}\)
are isomorphic. Their complete extensions
are exactly the sets of their respective
arguments whose negative supports are
contained in a common set \(\neg B\),
including the unit arguments for \(B\),
where
\[
 B=H\setminus\Gamma(\Gamma(H\setminus B)),
                  \qquad B\cap\Gamma(H\setminus B)=\varnothing.
\]
The isomorphism preserves conclusion sets,
grounded extensions, preferred extensions
and stable extensions. It also preserves
the property that the grounded extension
is the intersection of preferred extensions.
\end{hbtheorem}

\begin{hbdefinition}{def:p18-program-dependency}{The source predicate conditions}
For a finite source program, its finite predicate
graph has an edge from a head predicate to
each body predicate, signed positive or negative
according to the body literal. Walks may repeat
vertices. The program is hierarchical when
there is no nonempty closed walk; stratified
when no strongly connected component contains
a negative edge; strict when no ordered pair
admits both an even-negative and an odd-negative
nonempty walk; and call-consistent when there
is no odd-negative closed walk.
These are the conditions of Section~4.3.3.
\end{hbdefinition}

\begin{hbtheorem}{thm:p18-program-graph-transfer}{proved=annex:p18.program.graph.transfer}{Graph transfer at its valid argument domain}
For finite source programs:
\begin{enumerate}
\item Stratification makes \(\mathcal A_{\rm u}\) well-founded.
\item Hierarchy makes \(\mathcal A_{\rm c}\) well-founded.
\item Strictness makes both \(\mathcal A_{\rm u}\)
      and \(\mathcal A_{\rm c}\) uncontroversial.
\item Call-consistency makes both those graphs
      limited controversial.
\end{enumerate}
The semantic conclusions transfer back to
the literal support graph through the
preceding isomorphism. The graph assertions
themselves do not in general hold on
\(\mathcal A_\infty\): even the fact
\(p\leftarrow\) admits its argument
\((\{\neg p\},p)\), which attacks itself.
\end{hbtheorem}

\begin{hbtheorem}{thm:p18-program-coincidence}{proved=annex:p18.program.coincidence}{All resulting semantic coincidence and existence claims}
For finite source programs, allowing infinite grounding:
\begin{enumerate}
\item A stratified program has a unique stable
      model; its total literal interpretation
      is the well-founded model.
\item A hierarchical program has a unique
      Herbrand completion model, whose total
      literal interpretation is the Fitting model.
\item For a strict program and each ground
      literal \(l\), \(l\in\operatorname{WFM}(Q)\)
      exactly when \(l\) is true in every
      stable model. Preferred and stable
      extension semantics coincide.
\item For a call-consistent program the
      preferred extensions of each of the
      two source graphs are stable.
      Every maximal partial completion
      model is total.
\item Every call-consistent program has a
      stable model and a Herbrand completion
      model.
\end{enumerate}
These preserve the semantic conclusions
of Corollaries~59--61 with the explicit
support-reduction premise needed for
their graph route.
\end{hbtheorem}

\section{Effective domains and separating programs}

\begin{hbtheorem}{thm:p18-program-computability}{proved=annex:p18.program.computability}{Decidable clauses, semidecidable supports and the exact undecidability claim}
For a finite first-order source program,
membership of a supplied finite ground pair
in \(\mathcal A_{\rm c}\) is decidable.
Membership in \(\mathcal A_\infty\) is
semidecidable and is not uniformly decidable
in general. Uniform truth queries for
the well-founded model, stable semantics
and preferred extension semantics are
also undecidable in general, already
on definite programs with unique
stable/preferred conclusions.

For a finite explicitly grounded program
over finite \(H\), both argument graphs
are finite and can be constructed exactly.
All extension families can then be
computed by finite subset enumeration.
On infinite inputs a decidable argument
relation alone does not establish a
terminating semantics procedure.
\end{hbtheorem}

\begin{hbdefinition}{def:p18-program-examples}{Three small programs and the exact Appendix A program}
Take the displayed atoms as the entire
base in each propositional example:
\[
 Q_{\rm delay}=\{p\leftarrow\neg q\},\qquad
 Q_{\rm loop}=\{p\leftarrow p\},\qquad
 Q_{\rm fact}=\{p\leftarrow\}.
\]
The Appendix~A source program is
\[
\begin{array}{ll}
 r\leftarrow\neg p,&p\leftarrow\neg q(x),\\
 q(x)\leftarrow e(x),&q(x)\leftarrow\neg e(x),\\
 e(s(x))\leftarrow\neg e(x),&e(0)\leftarrow.
\end{array}
\]
Here \(e\) abbreviates the source predicate
\(\mathrm{even}\), and the ground terms
are \(s^n(0)\), \(n<\omega\).
\end{hbdefinition}

\begin{hbtheorem}{thm:p18-program-examples}{proved=annex:p18.program.examples}{Complete separating calculations}
For \(Q_{\rm delay}\), the first Fitting
iterate is \(\{\neg q\}\), whereas the
first clause-graph grounded conclusions
are \(\{\neg q,p\}\); both least fixed
points equal \(\{\neg q,p\}\).
For \(Q_{\rm loop}\), the well-founded
model is \(\{\neg p\}\) and the Fitting
model is empty. Its only stable program
model is empty, whereas its total
completion models are empty and \(\{p\}\).
For \(Q_{\rm fact}\), the literal
support graph contains the self-attacker
specified above, but both source
semantics have the single total
conclusion set \(\{p\}\).

For the Appendix~A program, both source
graphs are non-finitary. In the support
graph, writing \(e_n=e(s^n(0))\) and
\(q_n=q(s^n(0))\), the finite grounded
conclusion stages are \(I_0=\varnothing\) and
\[
 I_m=\{e_{2j}:j<m\}\cup\{q_n:n\leq2m-2\}
               \cup\{\neg e_{2j+1}:j<m-1\}\quad(m\geq1).
\]
At stage \(\omega\) their union lacks
\(r,\neg p\); both are added at
\(\omega+1\), which is fixed and total.
Thus this finite source program really
requires a transfinite ground-argument
stage.
\end{hbtheorem}

The default-theory correspondence of Section~4.1 is supplied in
\cref{thm:p18-default-stable-bijection}, with its separately proved
autoepistemic interface in
\cref{thm:p18-default-ael-weak,thm:p18-default-ael-prerequisite-free,thm:p18-default-ael-boundary}.
The game correspondence of Section~3 is supplied in
\cref{sec:p18-game-argumentation}, with the complete no-solution
construction in \cref{sec:p18-lucas-no-solution}.
The full Section~5 meta-interpreter, its guarded and sparse
specializations, and the stable, grounded and complete correspondences
are supplied in \cref{sec:p18-argumentation-meta-interpreters}. The present
chapter concerns Section~4.3 and its
required Appendices~A--C.
